% =================================================
% Учительская версия рабочего листа
% =================================================

\worksheettitle
  {Ломаная и многоугольник}
  {Учительская версия: ответы и методические комментарии}

\begin{teacherbox}[Цель листа]
  Сформировать различение ломаной и многоугольника; научить называть вершины,
  звенья и стороны, читать имя многоугольника по порядку обхода, проводить
  диагонали и называть фигуры, образованные дополнительными отрезками.
\end{teacherbox}

\begin{checkpointbox}[Ответ к входной разминке]
  Получился треугольник \(ABC\). Важно принять и названия \(BCA\), \(CAB\),
  \(ACB\), \(CBA\), \(BAC\): вершины идут подряд по границе.
\end{checkpointbox}

\section{Ломаная}

\openbrokenfigure

\begin{propertybox}[Заполненное определение]
  Ломаная состоит из последовательно соединённых \textbf{отрезков}. Эти
  отрезки называют \textbf{звеньями}, а их концы --- \textbf{вершинами}
  ломаной.
\end{propertybox}

\begin{practicebox}[Ответы к заданию 1]
  \begin{enumerate}[label=\alph*),itemsep=0.4em]
    \item Пять вершин.
    \item Четыре звена.
    \item \(AB\), \(BC\), \(CD\), \(DE\).
  \end{enumerate}
\end{practicebox}

\begin{teacherbox}[Методический акцент]
  Типичная ошибка --- считать число букв в названии числом звеньев. Для
  незамкнутой ломаной с \(k\) вершинами звеньев \(k-1\). Не стоит пока
  превращать это наблюдение в формулу: достаточно считать на рисунке.
\end{teacherbox}

\section{Замкнутая ломаная}

\studentrecognizefigure

\begin{practicebox}[Ответы к заданию 2]
  \begin{enumerate}[label=\alph*),itemsep=0.45em]
    \item Незамкнутые: 1 и 3.
    \item Замкнутые: 2 и 4.
    \item Дополнительная диагональ проведена на рисунке 4.
  \end{enumerate}
\end{practicebox}

\begin{teacherbox}[Важное различение]
  На рисунке 4 граница многоугольника остаётся замкнутой ломаной, а
  проведённая внутри диагональ не становится её звеном. Это подготавливает
  различение стороны и диагонали.
\end{teacherbox}

\section{Многоугольник}

\polygonpartsfigure

\begin{propertybox}[Ответ]
  Звенья границы называют \textbf{сторонами}, их концы ---
  \textbf{вершинами}. Число сторон \textbf{равно} числу вершин.
\end{propertybox}

\begin{practicebox}[Ответы к заданию 3]
  \begin{enumerate}[label=\alph*),itemsep=0.4em]
    \item Шесть вершин.
    \item Шесть сторон.
    \item Шестиугольник.
    \item \(AB\), \(BC\), \(CD\), \(DE\), \(EF\), \(FA\).
  \end{enumerate}
\end{practicebox}

\polygonnamesfigure

\begin{practicebox}[Ответы к заданию 4]
  Треугольник, пятиугольник, шестиугольник.
\end{practicebox}

\section{Название многоугольника}

\correctnamingfigure

\begin{thinkbox}[Ответ к заданию 5]
  Верны названия \(ABCDEF\), \(BCDEFA\), \(AFEDCB\). Запись \(ACDEFB\)
  неверна: после вершины \(A\) указана несоседняя вершина \(C\), то есть
  нарушен обход границы.
\end{thinkbox}

\begin{teacherbox}[Диагностика]
  Если ученик считает правильным любой порядок букв, попросите провести
  пальцем путь, заданный его записью. Переход через внутренность фигуры сразу
  показывает нарушение.
\end{teacherbox}

\clearpage
\section{Диагонали}

\diagonalfigure

\begin{propertybox}[Ответ]
  Диагональ соединяет две \textbf{несоседние} вершины. Из \(A\) проведены
  \(AC\), \(AD\), \(AE\).
\end{propertybox}

\begin{practicebox}[Ответ к заданию 6]
  \hexagondiagonalsanswer
  Всего три диагонали.
\end{practicebox}

\begin{teacherbox}[Почему \(n-3\)]
  Из выбранной вершины нельзя проводить диагональ к самой этой вершине и к
  двум соседним. Поэтому из \(n\) вершин остаётся \(n-3\). Формулу лучше
  выводить после нескольких рисунков, а не сообщать как правило для
  заучивания.
\end{teacherbox}

\section{Многоугольники внутри многоугольника}

\pentagonsplitfigure

\begin{practicebox}[Ответ к заданию 7]
  Диагональ \(AC\) образует треугольник \(ABC\) и четырёхугольник
  \(ACDE\). Допустимы циклические переименования и обратный порядок обхода,
  например \(BCA\) и \(EDCA\).
\end{practicebox}

\quadrilateralregionsfigure

\begin{thinkbox}[Ответ к заданию 8]
  После обозначения точки пересечения \(P\) получаются четыре области:
  \[
    \triangle ABP,\qquad \triangle BCP,\qquad \triangle APM,
    \qquad \text{четырёхугольник }PCDM.
  \]
\end{thinkbox}

\begin{teacherbox}[Методический акцент]
  Ученик может назвать более крупные многоугольники, составленные из
  нескольких областей, например \(ABC\). Формулировка «образовавшиеся
  части» означает именно минимальные области, на которые проведённые
  отрезки разделили исходную фигуру. Это стоит проговорить заранее.
\end{teacherbox}

\begin{checkpointbox}[Ответы к заданию 9]
  \begin{enumerate}[label=\alph*),itemsep=0.4em]
    \item \(+\)
    \item \(-\): соседние вершины соединяет сторона.
    \item \(+\)
    \item \(-\): вершины перечисляют подряд по границе.
    \item \(+\), так как \(7-3=4\).
  \end{enumerate}
\end{checkpointbox}

\section{Типичные ошибки}

\begin{center}
  \renewcommand{\arraystretch}{1.42}
  \begin{tabularx}{\textwidth}{
    >{\raggedright\arraybackslash}p{0.29\textwidth}
    >{\raggedright\arraybackslash}p{0.32\textwidth}
    >{\raggedright\arraybackslash}X}
    \toprule
    \rowcolor{TableHeader}
    Ошибка & Что не понято & Наводящий вопрос \\
    \midrule
    Число звеньев равно числу вершин у незамкнутой ломаной
      & Не различены точки и переходы между ними
      & Сколько отрезков реально проведено? \\
    \addlinespace[0.2em]
    В названии многоугольника буквы переставлены произвольно
      & Имя должно задавать обход границы
      & Можно ли пройти по сторонам в таком порядке? \\
    \addlinespace[0.2em]
    Сторона названа диагональю
      & Не проверено соседство вершин
      & Эти вершины соседние или несоседние? \\
    \addlinespace[0.2em]
    При перечислении областей пропущена точка пересечения
      & Граница части прочитана не полностью
      & Через какие отмеченные точки проходит граница? \\
    \bottomrule
  \end{tabularx}
\end{center}

\begin{teacherbox}[Итоговая формулировка ученика]
  «Многоугольник называют, перечисляя вершины подряд по границе. Стороны
  соединяют соседние вершины, а диагонали --- несоседние».
\end{teacherbox}
